\documentclass[11pt]{article}
\usepackage[letterpaper,margin=0.83in,headheight=16pt]{geometry}
\usepackage[T1]{fontenc}
\usepackage[utf8]{inputenc}
\usepackage{lmodern}
\usepackage{microtype}
\usepackage{amsmath,amssymb}
\usepackage{xcolor}
\usepackage{hyperref}
\usepackage{enumitem}
\usepackage{booktabs,longtable,tabularx,array}
\usepackage{fancyhdr}
\usepackage{titlesec}
\usepackage{listings}
\usepackage{tikz}
\usepackage{float}
\definecolor{suiteblue}{HTML}{17365D}
\definecolor{suitelight}{HTML}{EAF1F8}
\definecolor{suitegray}{HTML}{556070}
\hypersetup{colorlinks=true,linkcolor=suiteblue,urlcolor=suiteblue,breaklinks=true}
\titleformat{\section}{\Large\bfseries\color{suiteblue}}{\thesection}{0.55em}{}
\titleformat{\subsection}{\large\bfseries\color{suiteblue}}{\thesubsection}{0.55em}{}
\setlist[itemize]{itemsep=0.25em,topsep=0.35em,leftmargin=1.4em}
\setlist[enumerate]{itemsep=0.25em,topsep=0.35em,leftmargin=1.6em}
\setlength{\parindent}{0pt}
\setlength{\parskip}{0.47em}
\emergencystretch=2em
\lstset{basicstyle=\ttfamily\footnotesize,breaklines=true,frame=single,columns=fullflexible,showstringspaces=false}
\pagestyle{fancy}\fancyhf{}
\lhead{Game Design \& Engineering}\rhead{Workbook Template}\cfoot{\thepage}
\begin{document}
{\huge\bfseries\color{suiteblue} Game Design Architecture Workbook -- Template}\par
\vspace{0.15in}
{\normalsize\color{suitegray} Fictional example annotations plus reusable analytical prompts}\par
\vspace{0.15in}\hrule\vspace{0.15in}

\section{Record control}
Project: [Project]. Revision: [Revision]. Practitioner: [Name]. Evidence repository: [URI].


\section{Evidence discipline}
For each exercise, distinguish an assumption, result established by computation, result established by player observation, and an interpretation. The example entries are instructional and are not observations from a real game.


\section{Boundary mapping}
Exercise prompt: Allocate input validation, simulation state, presentation, persistence, and analytics to distinct responsibilities.

Expected deliverable: Component graph and owner table.


\subsection{Hypothesis}
Illustrative example: A small design change affects a specified, observable player action.

Project-specific response: [Document the evidence and rationale.]


\subsection{Assumptions}
Illustrative example: The simulated population and state constraints are explicitly defined.

Project-specific response: [Document the evidence and rationale.]


\subsection{Model or procedure}
Illustrative example: Use a repeatable implementation and record input parameters.

Project-specific response: [Document the evidence and rationale.]


\subsection{Results and limitations}
Illustrative example: Quantitative result is supported by a run log; enjoyment remains unverified.

Project-specific response: [Document the evidence and rationale.]


\subsection{Decision and next test}
Illustrative example: Keep, revise, or retire the feature after a new independent study.

Project-specific response: [Document the evidence and rationale.]


\section{Quality scenario}
Exercise prompt: Given a target of responding promptly to a controller action, define stimulus, response, environmental conditions, and a measurable test budget.

Expected deliverable: Quality-attribute scenario and test plan.


\subsection{Hypothesis}
Illustrative example: A small design change affects a specified, observable player action.

Project-specific response: [Document the evidence and rationale.]


\subsection{Assumptions}
Illustrative example: The simulated population and state constraints are explicitly defined.

Project-specific response: [Document the evidence and rationale.]


\subsection{Model or procedure}
Illustrative example: Use a repeatable implementation and record input parameters.

Project-specific response: [Document the evidence and rationale.]


\subsection{Results and limitations}
Illustrative example: Quantitative result is supported by a run log; enjoyment remains unverified.

Project-specific response: [Document the evidence and rationale.]


\subsection{Decision and next test}
Illustrative example: Keep, revise, or retire the feature after a new independent study.

Project-specific response: [Document the evidence and rationale.]


\section{Trust boundary}
Exercise prompt: Model hostile saved-game input, forged client events, and opt-in telemetry.

Expected deliverable: Trust-boundary diagram and control table.


\subsection{Hypothesis}
Illustrative example: A small design change affects a specified, observable player action.

Project-specific response: [Document the evidence and rationale.]


\subsection{Assumptions}
Illustrative example: The simulated population and state constraints are explicitly defined.

Project-specific response: [Document the evidence and rationale.]


\subsection{Model or procedure}
Illustrative example: Use a repeatable implementation and record input parameters.

Project-specific response: [Document the evidence and rationale.]


\subsection{Results and limitations}
Illustrative example: Quantitative result is supported by a run log; enjoyment remains unverified.

Project-specific response: [Document the evidence and rationale.]


\subsection{Decision and next test}
Illustrative example: Keep, revise, or retire the feature after a new independent study.

Project-specific response: [Document the evidence and rationale.]


\section{Change analysis}
Exercise prompt: Change resource cost from one to two. Identify affected tests, tutorials, balancing models, and save migration requirements.

Expected deliverable: Impact traceability links and change approval.


\subsection{Hypothesis}
Illustrative example: A small design change affects a specified, observable player action.

Project-specific response: [Document the evidence and rationale.]


\subsection{Assumptions}
Illustrative example: The simulated population and state constraints are explicitly defined.

Project-specific response: [Document the evidence and rationale.]


\subsection{Model or procedure}
Illustrative example: Use a repeatable implementation and record input parameters.

Project-specific response: [Document the evidence and rationale.]


\subsection{Results and limitations}
Illustrative example: Quantitative result is supported by a run log; enjoyment remains unverified.

Project-specific response: [Document the evidence and rationale.]


\subsection{Decision and next test}
Illustrative example: Keep, revise, or retire the feature after a new independent study.

Project-specific response: [Document the evidence and rationale.]


\section{Release gate}
Exercise prompt: Compare functional evidence and playtest evidence; classify which claims can and cannot be supported by each.

Expected deliverable: Release readiness memo and open risk register.


\subsection{Hypothesis}
Illustrative example: A small design change affects a specified, observable player action.

Project-specific response: [Document the evidence and rationale.]


\subsection{Assumptions}
Illustrative example: The simulated population and state constraints are explicitly defined.

Project-specific response: [Document the evidence and rationale.]


\subsection{Model or procedure}
Illustrative example: Use a repeatable implementation and record input parameters.

Project-specific response: [Document the evidence and rationale.]


\subsection{Results and limitations}
Illustrative example: Quantitative result is supported by a run log; enjoyment remains unverified.

Project-specific response: [Document the evidence and rationale.]


\subsection{Decision and next test}
Illustrative example: Keep, revise, or retire the feature after a new independent study.

Project-specific response: [Document the evidence and rationale.]


\section{Review decision}
Reviewer: [Name]. Findings: [Finding IDs]. Status: [Draft, revision needed, or accepted].


\end{document}
